\begin{frame}{Quokka 原流程（未改演算法）}
  \vspace{-0.2em}
  \centering
  \begin{tikzpicture}[
  node distance=0.55cm and 0.9cm,
  box/.style={draw=Accent, rounded corners=2pt, minimum height=0.75cm,
    minimum width=2.1cm, align=center, font=\TikzCJKFont},
  arr/.style={-{Stealth[length=2mm]}, thick, Accent}
]
  \node[box] (p) {程式\\迴圈點};
  \node[box, right=of p] (l) {LLM\\assume};
  \node[box, right=of l] (a) {assume\\驗證};
  \node[box, right=of a] (b) {assert\\驗證};
  \node[box, right=of b] (m) {\#Corr\\\#Ext@T};
  \draw[arr] (p) -- (l) -- (a) -- (b) -- (m);
\end{tikzpicture}

  \vspace{0.5em}
  {\footnotesize\textcolor{SeriesPaper}{soundness 來自 UAutomizer · 不信 LLM 原文}\par}
  \note{
    \textbf{約 45 秒}

    Quokka 的核心沒改：LLM 依 prompt 產生要在某一行插入的 assume 不變式；對每個 sample 做 UAutomizer 兩階段驗證，先 assume 再 assert。只要有一個 sample 的 assert 是 TRUE，這題就算進 \#Corr。指標是和 timing\_uautomizer 的 baseline 比，算 \#Ext 在 T 秒內多解了幾題、以及 PAR。重點是：正確性來自 verifier，不是相信模型原文。下一頁會說：一次 LLM 呼叫只產一條 invariant，且模型從多個 loop 插入點裡自選一個。
  }
\end{frame}
